\chapter*{Abstract}
\addcontentsline{toc}{chapter}{Abstract}

Steel-truss sizing under IS~800:2007 is classically solved with
evolutionary metaheuristics whose inner loops spend almost all
wall-clock time inside a finite-element solver. This thesis develops a
reproducible Python framework that combines classical evolutionary
optimisation (GA, PSO, NSGA-II) with three modern AI layers --- a neural
surrogate, a reinforcement-learning agent, and an AI design-agent
warm-start --- and validates it on four canonical benchmarks: the
10-bar planar, 25-bar spatial, 72-bar spatial tower, and 200-bar stepped
tower.

The framework reproduces published optima on 10-bar to
\num{5061.05}~lb (\num{0.004}~\% error) and 25-bar to \num{545.26}~lb
(\num{0.008}~\% error). A three-hidden-layer MLP surrogate trained on
\num{10000} Latin-hypercube samples reaches $R^2=\num{0.9993}$ on weight
and gives a $\num{132}\times$ wall-clock speedup, with
surrogate-in-loop optima matching FEM to $\num{0.09}\,\%$. On the 72-bar
tower the hard-constraint solvers converge with cross-seed spread under
$\num{1}\,\%$ to a feasible optimum near $\num{549}$~lb; we flag the gap
versus the soft-penalty $\num{380}$~lb literature value as an observation
pending an independent FEM cross-check. An agent warm-start reduces
generations-to-convergence on 10-bar by \num{36.8}~\%
($p=\num{0.0046}$), with the effect diminishing on larger benchmarks
where no comparable architectural insight is available. Every reported
optimum passes the IS~800:2007 slenderness, tension, compression, and
serviceability checks, and the full pipeline --- shipped at tag
\texttt{phase-10-complete} --- reproduces every headline number in under
ten seconds on a laptop CPU, offline, with all agent responses cached.

\clearpage
